%% ============================================================================
%% ch4_market.tex -- Chapter 4: חקר השוק והמוצרים הקיימים
%% budget: body pages 17-22 (6 pages)
%% Market figures = team estimates (no fabricated citations). Charts read
%% figures/data/radar.csv and scatter.csv (scripts/generate_beauty_benchmarks.py).
%% ============================================================================
\chapter{חקר השוק והמוצרים הקיימים}
\label{ch:market}

\begin{managerbox}
דיוק טכני אינו שווה דבר אם אין מי שישלם עליו, או אם מתחרה כבר מוכר אותו. הפרק מראה שקיים שוק בהיקף משמעותי, שהמוצרים הקיימים פותרים כל אחד חלק אחר של הבעיה, ושמנוע החלפת הערך מכסה את עלות החישוב. גם הסטייליסטית האנושית והזמנת שתי מידות נבחנות כאן כמתחרות.
\end{managerbox}

\begin{warningbox}[מדיניות מספרי השוק]
מספרי הפרק הם משני סוגים. היקף השוק הכולל ועובדות על המתחרים לקוחים ממקורות שנפתחו ונבדקו, והמקור מצוין לצדם; תקצירים של חברות מחקר הם מקור משני, ולכן הם משמשים לסדר גודל בלבד. ה-\en{SAM}, ה-\en{SOM}, טבלת ההשוואה וכלכלת היחידה הם הנחות של הצוות, מסומנות ככאלה ונבנות מלמטה למעלה, כך שכל אחת ניתנת לבדיקה בפיילוט.
\end{warningbox}

\section{היקף השוק בשלוש שכבות}
\label{sec:tam}

ניתוח השוק נבנה במודל השכבות (איור~\ref{fig:tam}). \en{TAM} הוא השוק העולמי שבו הפתרון פועל: שוק חדרי המדידה הווירטואליים נאמד ב-\num{\$6.9B} בשנת \num{2025}, עם \en{CAGR} של \num{17.7\%} עד \num{2034} \cite{fbi2026vfr}, ושוק הבינה המלאכותית בביוטי ב-\num{\$3.5B} בשנת \num{2023}, עם \en{CAGR} של \num{19.1\%} \cite{kings2024aibeauty}. יחד זהו סדר גודל של \num{\$10B}; הוא אינו משמש לשום החלטה בדוח. \en{SAM} הוא הפלח שאליו פונה הפתרון באירופה, בישראל ובארצות הברית: קמעונאיות שהכנסתן המקוונת מעל \num{\$50M} ומשתמשות של אפליקציות ביוטי ואופנה, בכ-\num{\$4.7B}. \en{SOM} הוא הנתח המעשי בשנה השלישית, בכ-\num{\$45.0M}. שני האחרונים נבנים מלמטה בטבלה~\ref{tab:som}, מהנחות שכל אחת מהן ניתנת לבדיקה בנפרד.

\begin{reportfigure}{שלוש שכבות השוק של \en{AURAM} (\en{TAM} ממקורות; \en{SAM} ו-\en{SOM} אומדני הצוות לשנה השלישית)}{fig:tam}
\begin{tikzpicture}[x=0.88cm,y=0.88cm,
  tag/.style={font=\small\bfseries,text=#1},
  lead/.style={draw=brandprimary!70,line width=0.7pt},
  dot/.style={circle,fill=brandprimary!70,inner sep=1.3pt},
  lab/.style={anchor=west,align=left,font=\footnotesize,text=brandprimary}]
  % nested, bottom-aligned layers in three clearly different solid colours; no text over boundaries
  \fill[brandprimary] (0,2.4) circle (2.4);
  \fill[brandsecondary] (0,1.55) circle (1.55);
  \fill[brandhighlight] (0,0.75) circle (0.75);
  \draw[brandaccent,line width=1pt] (0,2.4) circle (2.4) (0,1.55) circle (1.55) (0,0.75) circle (0.75);
  \node[tag=brandcream] at (0,4.3) {TAM};
  \node[tag=brandcream] at (0,2.68) {SAM};
  \node[tag=white] at (0,0.75) {SOM};
  % leader lines to a label column on the right
  \draw[lead] (1.5,3.9) -- (3.0,3.9); \node[dot] at (1.5,3.9) {};
  \draw[lead] (0.9,2.3) -- (3.0,2.3); \node[dot] at (0.9,2.3) {};
  \draw[lead] (0.42,0.75) -- (3.0,0.75); \node[dot] at (0.42,0.75) {};
  \node[lab] at (3.1,3.9) {\textbf{TAM $\approx$ \$10B (2023--2025)}\\[-1pt]{\scriptsize\heb{מדידה וירטואלית ובינה מלאכותית בביוטי; צמיחה שנתית של \num{17\%--19\%}}}};
  \node[lab] at (3.1,2.3) {\textbf{SAM $\approx$ \$4.7B}\\[-1pt]{\scriptsize\heb{אירופה, ישראל וארצות הברית: \num{3,000} קמעונאיות ו-\num{150M} משתמשות אפליקציות}}};
  \node[lab] at (3.1,0.75) {\textbf{SOM $\approx$ \$45.0M} \heb{\textbf{(שנה שלישית)}}\\[-1pt]{\scriptsize\heb{\num{12} קמעונאיות ו-\num{1.6M} משתמשות; כ-\num{1.0\%} מה-\en{SAM}}}};
\end{tikzpicture}
\end{reportfigure}

ה-\en{SOM} הוא כ-\num{1.0\%} מה-\en{SAM}, נתח צנוע ביחס לשוק. ההכנסה החודשית הממוצעת למשתמשת (\en{ARPU}), \num{\$1.59}, היא הפלט של מנוע החלפת הערך בסעיף~\ref{sec:unit}: עמלות על המלצות \num{\$0.30}, הכנסות סביב הדמיות \num{\$1.09} (דמי חסות, מיקרו-אינטראקציות ותוספת ההמרה שההדמיה יוצרת) ומנוי נטו \num{\$0.20}.

\begin{reporttable}[\footnotesize]{בניית ה-\en{SAM} וה-\en{SOM} מלמטה למעלה (אומדן הצוות, שנה שלישית)}{tab:som}{R P{6.6cm} P{2.4cm}}
\headrow \hc{שכבה וערוץ} & \hc{הנחת חישוב} & \hc{הכנסה שנתית} \\ \midrule
\en{SAM}, קמעונאיות & קמעונאיות $\times$ רישוי: \en{3,000 $\times$ \$0.6M} & \num{\$1.80B} \\
\en{SAM}, משתמשות & פעילות $\times$ הכנסה $\times$ חודשים: \en{150M $\times$ \$1.59 $\times$ 12} & \num{\$2.87B} \\
\en{SOM}, קמעונאיות & שותפות $\times$ הזמנות $\times$ עמלה: \en{12 $\times$ 2M $\times$ \$0.60} & \num{\$14.4M} \\
\en{SOM}, משתמשות & פעילות $\times$ הכנסה $\times$ חודשים: \en{1.6M $\times$ \$1.59 $\times$ 12} & \num{\$30.6M} \\
\hlrow \textbf{\en{SOM} סה"כ} & & \textbf{\num{\$45.0M}} \\
\end{reporttable}

ה-\en{SOM} נשען על שני ערוצים, \num{\$14.4M} מקמעונאיות ו-\num{\$30.6M} ממשתמשות; שניהם תלויים בהמרות אצל אותם קמעונאים, ולכן הפיזור חלקי. ב-\en{SAM} נלקח רישוי ממוצע של \num{\$0.6M} לקמעונאית בינונית, וב-\en{SOM} נבחרו \num{12} קמעונאיות מהשורה הראשונה, שהתמחור שלהן לפי הזמנה מגיע ל-\num{\$1.2M}. ההנחה הרגישה ביותר היא \num{1.6} מיליון משתמשות פעילות (\en{RSK-OPS-06}, פרק~\ref{ch:risk}).

\section{סקירה טכנית של המוצרים הקיימים}
\label{sec:competitors}

\textbf{מוצרי מציאות רבודה.} \en{ModiFace}, שנרכשה בידי \en{L'Oréal} ב-\num{2018} \cite{loreal2018modiface}, ו-\en{YouCam Makeup} של \en{Perfect Corp}, שמדווחת על \num{859} מותגים לקוחות \cite{perfect2026results}, מאתרות נקודות ציון בפנים ומציירות שכבת איפור בזמן אמת, בעשרות פריימים לשנייה. נקודת התורפה משותפת: השכבה מצוירת מעל תמונה שכבר אוזנה בידי המצלמה, בלי כיול צבע ביומטרי, ואין טיפול במידות גוף או בלבוש. \textbf{הלבשה וירטואלית.} \en{Zeekit}, שנרכשה בידי \en{Walmart} ב-\num{2021} \cite{walmart2021zeekit}, מפלחת את הגוף ומעוותת אליו את תמונת הבגד; היא אינה מטפלת בצבע עור, באיפור או בציפורניים. \textbf{סטיילינג אנושי.} \en{Stitch Fix} משלבת אלגוריתמי המלצה עם סטייליסטיות ושולחת ארגז פריטים תמורת דמי סטיילינג של \num{\$20} \cite{stitchfix2026tenk}: שיפוט אנושי ומשוב רכישה אמיתי, במחיר של ימי המתנה ועלות כוח אדם לכל לקוחה. \textbf{החלופה הידנית,} יועצת בחנות או הזמנת שתי מידות והחזרת אחת, היא המתחרה בפועל: מדויקת יחסית, ויקרה בזמן ובלוגיסטיקה.

טבלה~\ref{tab:compare} מעמתת את המוצרים עם מאמר הדגל ועם הפתרון המתוכנן, לפי שמונה פרמטרים הנדסיים.

\begin{reporttable}[\footnotesize]{השוואה רב-פרמטרית: מוצרים מסחריים, מאמר הדגל והפתרון המוצע (הערכת הצוות)}{tab:compare}{P{2.3cm} C C C C C C}
\headrow \hc{פרמטר} & \hc{\en{ModiFace}} & \hc{\en{YouCam}} & \hc{\en{Zeekit}} & \hc{\en{Stitch}\newline\en{Fix}} & \hc{\en{IDM-}\newline\en{VTON}} & \hc{\en{AURAM}} \\ \midrule
טכנולוגיה & \en{AR} פנים & \en{AR} פנים ושיער & עיוות בגד & טבלאי ואנושי & דיפוזיה & היברידי \\
כיול צבע & חלקי & חלקי & אין & אין & אין & עוגן עין \\
מידות גוף & אין & אין & צללית & שאלון & אין & \en{SMPL-X} וגובה \\
פוטו-ריאליזם & בינוני & בינוני & בינוני-גבוה & אין & גבוה מאוד & גבוה \\
זמן תגובה & \num{< 0.1 s} & \num{< 0.1 s} & שניות & ימים & \num{> 8 s} & \num{1.8 s}\newline\num{4.5 s} \\
עלות לסשן & \num{\$0.004} & \num{\$0.003} & \num{\$0.03} & \num{\$8} & \num{\$0.09} & \num{\$0.025} \\
קטגוריות & איפור & איפור, שיער & לבוש & לבוש & לבוש & כל הארבע \\
מודל הכנסה & רישוי & מנוי, \en{SDK} & קמעונאי & עמלת סטיילינג & מחקר & חסות ושותפים \\
\end{reporttable}

כל מוצר מצטיין בשורה אחת או שתיים: מוצרי המציאות הרבודה זולים ומהירים אך אינם מכיילים ואינם מודדים, ומאמר הדגל מגיע לפוטו-ריאליזם הגבוה ביותר במחיר הגבוה ביותר. השורות "כיול צבע" ו"מידות גוף" ריקות כמעט לגמרי, כמו העמודות המתאימות במטריצת הספרות (טבלה~\ref{tab:litmatrix}): הפער שזוהה באקדמיה קיים גם בשוק.

\section{מיצוב: פרופיל ביצועים ודיוק מול עלות}
\label{sec:graphs}

איור~\ref{fig:radar} מציג את פרופיל הביצועים של כל פתרון על פני שישה צירים מנורמלים לסולם \num{0--5}. הערכים הם הערכה טכנית של הצוות, על בסיס הניתוח שלמעלה, ולא מדידה.

\begin{reportfigure}{פרופיל ביצועים: \en{AURAM} אינה הגבוהה בשום ציר, אבל היחידה בלי ציר חלש (הערכת הצוות, סולם \num{0--5})}{fig:radar}
\begin{tikzpicture}
\begin{polaraxis}[width=6.3cm, ymin=0, ymax=5, ytick={1,2,3,4,5}, yticklabel style={font=\scriptsize},
  xtick={0,60,120,180,240,300},
  xticklabels={\heb{כיול צבע},\heb{התאמת גזרה},\heb{פוטו-ריאליזם},\heb{מהירות},\heb{כיסוי קטגוריות},\heb{רווח לחישוב}},
  xticklabel style={font=\scriptsize}, grid style={black!15},
  legend style={at={(1.32,0.5)},anchor=west,font=\scriptsize,draw=black!20}]
  \addplot[brandgrey,thick,mark=none] table[x=angle,y=modiface]{figures/data/radar.csv};
  \addplot[brandgrey!60!black,thick,dashed,mark=none] table[x=angle,y=youcam]{figures/data/radar.csv};
  \addplot[brandsecondary,thick,dotted,mark=none] table[x=angle,y=zeekit]{figures/data/radar.csv};
  \addplot[brandprimary!50,thick,dashdotted,mark=none] table[x=angle,y=stitchfix]{figures/data/radar.csv};
  \addplot[brandaccent!90!black,very thick,mark=none] table[x=angle,y=flagship]{figures/data/radar.csv};
  \addplot[brandhighlight,ultra thick,fill=brandhighlight,fill opacity=0.15,mark=*,mark size=1.2pt] table[x=angle,y=auratwin]{figures/data/radar.csv};
  \legend{ModiFace, YouCam, Zeekit, Stitch Fix, {\heb{\en{IDM-VTON} (מאמר הדגל)}}, {\heb{\en{AURAM} (יעד)}}}
\end{polaraxis}
\end{tikzpicture}
\end{reportfigure}

לכל מתחרה יש "קוץ" בכיוון אחד: מוצרי המציאות הרבודה במהירות, מאמר הדגל בפוטו-ריאליזם, ו-\en{Stitch Fix} בהתאמת הגזרה. הפרופיל המתוכנן של \en{AURAM} אינו הגבוה ביותר בשום ציר, אבל הוא היחיד שאין בו ציר נמוך מ-\num{4}. ערכי המתחרים הם הערכת הצוות וערכי \en{AURAM} הם יעד: הרדאר אינו הוכחה לעליונות אלא תרגום חזותי של יעדי המפרט, ובאבן דרך ב' הקו המתוכנן יוחלף בערכים שנמדדו.

איור~\ref{fig:scatter} מעמת את מדד דיוק ההתאמה האישית מול העלות לסשן בסולם לוגריתמי. ערכי הדיוק של המתחרים הם הערכה להמחשה ולא מדידה, ולכן לא הותאם להם קו מגמה.

\begin{reportfigure}{דיוק מול עלות לסשן: היעד של \en{AURAM} ביחס לפתרונות הקיימים (ערכים להמחשה, הערכת הצוות)}{fig:scatter}
\begin{tikzpicture}
\begin{semilogxaxis}[width=11.4cm, height=4.6cm, xmin=0.0015, xmax=25, ymin=45, ymax=95,
  xlabel={\heb{עלות לסשן סטיילינג (דולר, סולם לוגריתמי)}}, ylabel={\heb{מדד דיוק ההתאמה האישית (\num{0--100})}},
  label style={font=\small}, tick label style={font=\small}, grid=major, grid style={black!10},
  log ticks with fixed point, xtick={0.003,0.01,0.03,0.1,0.3,1,3,10}]
  \addplot[only marks, mark=*, mark size=3pt, brandsecondary,
    nodes near coords, point meta=explicit symbolic,
    every node near coord/.append style={font=\scriptsize,text=black!75,anchor=west,xshift=2pt}]
    table[x=cost,y=acc,meta=name,col sep=comma]{figures/data/scatter.csv};
  \addplot[only marks, mark=star, mark size=6pt, brandhighlight, very thick] coordinates {(0.025,88)};
\end{semilogxaxis}
\end{tikzpicture}
\end{reportfigure}

בין הפתרונות הקיימים, \en{Stitch Fix} מגיעה לדיוק הגבוה ביותר, בעלות גבוהה פי כמה מאות מזו של מוצרי המציאות הרבודה. היעד של \en{AURAM} הוא \num{88}, בעלות נמוכה פי \num{3.6} מזו של מאמר הדגל. סשן ממוצע כולל המלצה (\num{\$0.002}) וכ-\num{1.5} הדמיות בעלות כוללת של \num{\$0.015} כל אחת (סעיף~\ref{sec:unit}). המיקום הוא יעד תכן ולא מדידה, והוא ייבחן מול \en{REQ-VIS-01}, \en{REQ-BOD-03} ו-\en{REQ-BIZ-12}.

\section{מנוע החלפת הערך: שלוש שכבות במקום מודעה}
\label{sec:unit}

הדמיה גנרטיבית היא הרכיב היקר במערכת, אך מודעה של \num{15} שניות שאי אפשר לדלג עליה היא מלכודת חיכוך: לפי הנחת הצוות היא מאבדת כ-\num{80\%} מהמשתמשות ברגע ההחלטה. לכן השער מתוכנן כשלוש שכבות. \textbf{שכבה \num{1}, מיידית וחינמית:} התאמת גוון בקצה וצללית בגד פרמטרית מעל \en{SMPL-X}, בפחות מ-\num{300} אלפיות שנייה ובלי עלות \en{GPU}. \textbf{שכבה \num{2}, הדמיה ממומנת:} כשהפריט שייך לקמעונאי שותף, הוא משלם \num{\$0.02} להדמיה כיחידת מעורבות בעלת כוונת קנייה גבוהה. \textbf{שכבה \num{3}, החלפת ערך:} מיקרו-אינטראקציה ממותגת של \num{3} שניות, קרדיטים על ביקורת או העלאת ארון, או מנוי \en{VIP} של \num{\$7.99} לחודש. התרומה החודשית של משתמשת פעילה היא:
\begin{equation}
\begin{gathered}
n_g = n\,\big[s + (1-s)\,e^{-\beta t}\big], \qquad \frac{\text{LTV}}{\text{CAC}} = \frac{M}{\lambda\,\text{CAC}}, \\
M = R_{\text{aff}} + R_{\text{sub}} + n\,s\,f + n\,(1-s)\,e^{-\beta t}\,\frac{\text{eCPM}}{1000} \\
\qquad + \, n_g\big(\Delta p \cdot \text{AOV} \cdot r_{\text{aff}} - C_{\text{GPU}}\big) - C_{\text{rec}},
\end{gathered}
\label{eq:unit}
\end{equation}
כאשר $n$ הוא מספר בקשות ההדמיה בחודש, $s$ חלקן על פריטים ממומנים, $f$ דמי החסות, $\beta$ קצב הנטישה לשנייה של חיכוך ו-$t$ משך האינטראקציה, $n_g$ מספר ההדמיות שנוצרות, $\text{eCPM}$ ההכנסה לאלף מיקרו-אינטראקציות, $\Delta p$ תוספת הסתברות הרכישה אחרי הדמיה, $\text{AOV}$ ערך ההזמנה, $r_{\text{aff}}$ שיעור העמלה, $C_{\text{GPU}}$ העלות הכוללת של הדמיה, $C_{\text{rec}}$ עלות ההמלצות, $R_{\text{aff}}$ ו-$R_{\text{sub}}$ ההכנסה מעמלות ההמלצה ומהמנוי, $M$ התרומה החודשית, $\lambda$ שיעור הנטישה החודשי, $\text{LTV}$ ערך חיי הלקוחה ו-$\text{CAC}$ עלות הגיוס. הפרמטר $\beta = \ln 5/15$ נגזר מהנחת אובדן של \num{80\%} במודעה של \num{15} שניות, ולכן אינטראקציה של \num{3} שניות שומרת על \num{72\%} מהמשתמשות, מול \num{20\%} בלבד במודעה. טבלה~\ref{tab:unit} מציבה את הנחות הצוות.

\begin{reporttable}[\footnotesize]{מנוע החלפת הערך: הנחות הצוות ותוצאות לחודש של משתמשת פעילה}{tab:unit}{R P{4.6cm} P{4.5cm}}
\headrow \hc{רכיב} & \hc{הנחה} & \hc{ערך} \\ \midrule
בקשות הדמיה ונוצרות & \en{$n = 10.2$}, \en{$s = 0.6$}, \en{$t = 3$ s} & \num{$n_g = 9.08$} \\
עלות \en{GPU} כוללת & התחלה קרה, תור, מרווח \num{20\%} & \num{\$0.008--0.025} (בסיס \num{\$0.015}) \\
תוספת ההמרה בהדמיה & \en{4\% $\to$ 6\%}, סל \en{\$65}, עמלה \en{8\%} & \num{\$0.104} להדמיה \\
הכנסה חודשית & עמלות, הדמיות, מנוי & \num{\$1.59} \\
\hlrow תרומה חודשית $M$ & בניכוי \num{\$0.148} עלויות חישוב & \num{\$1.45} (\num{91\%}) \\
\hlrow \en{LTV/CAC} & \en{$\lambda = 12\%$}, \en{CAC = \$2.00} & \num{6.0} \\
\end{reporttable}

\begin{reporttable*}[\footnotesize]{\en{LTV/CAC} רגיש לנטישה החודשית, וכמעט אדיש לעלות ההדמיה}{tab:ltv}{R Q{3.0cm} Q{3.0cm} Q{3.0cm}}
\headrow \hc{עלות \en{GPU} להדמיה} & \hc{נטישה \num{8\%}} & \hc{נטישה \num{12\%}} & \hc{נטישה \num{20\%}} \\ \midrule
\num{\$0.008} & \num{9.43} & \num{6.29} & \num{3.77} \\
\num{\$0.015} & \num{9.03} & \num{6.02} & \num{3.61} \\
\num{\$0.025} & \num{8.47} & \num{5.64} & \num{3.39} \\
\end{reporttable*}

מהטבלאות עולה שסף היעד של הצוות, \num{$\text{LTV/CAC} > 3.8$}, מתקיים בכל טווח עלויות ה-\en{GPU} כל עוד הנטישה החודשית אינה עולה על \num{12\%}, ונכשל בנטישה של \num{20\%}. ההכנסה הצפויה מהדמיה אחת היא \num{\$0.120}, פי \num{8.0} מהעלות הבסיסית ופי \num{4.8} מהעלות הגבוהה, ולכן \en{REQ-BIZ-12} עומדת גם בתרחיש היקר.

\textbf{\num{6.0} הוא נקודה בתוך התפלגות.} כל אחת מ-\num{12} הנחות הצוות קיבלה טווח אי-ודאות במקום ערך יחיד, ו-\num{200,000} הגרלות מונטה קרלו עברו דרך משוואה~\eqref{eq:unit} (הטווחים הם הערכת הצוות ומתועדים בקוד, \en{scripts/uncertainty\_voi.py}). החציון הוא \num{6.1}, אבל \num{90\%} מההגרלות נופלות בין \num{2.2} ל-\num{18.3}, והסיכוי לעבור את סף \num{3.8} הוא \num{78\%}, לא ודאות (איור~\ref{fig:voi}, משמאל).

\begin{reportfigure}{\en{LTV/CAC} הוא התפלגות, וערך המידע מדרג את המדידות: ניסוי ההמרה אצל השותף ראשון, עלות ה-\en{GPU} אחרונה}{fig:voi}
\begin{tikzpicture}
\begin{groupplot}[group style={group size=2 by 1,horizontal sep=3.7cm}, height=3.6cm,
  tick label style={font=\scriptsize}, label style={font=\scriptsize}]
\nextgroupplot[width=5.9cm, ybar, bar width=2.3pt, xmin=0, xmax=20, ymin=0, ymax=6.2,
  xlabel={\en{LTV/CAC}}, ylabel={\heb{אחוז מההגרלות}}, ymajorgrids, grid style={black!10}]
  \addplot[fill=brandsecondary!70,draw=none] table[x=ratio,y=pct]{figures/data/ltv_cac_distribution.csv};
  \draw[brandwarn,dashed,thick] (axis cs:3.8,0) -- (axis cs:3.8,6.2);
  \draw[brandhighlight,thick] (axis cs:6.1,0) -- (axis cs:6.1,5.2);
  \node[font=\scriptsize,text=brandwarn,anchor=south east] at (axis cs:19.6,4.9) {\heb{סף היעד \num{3.8}}};
  \node[font=\scriptsize,text=brandhighlight,anchor=south east] at (axis cs:19.6,3.9) {\heb{חציון \num{6.1}}};
\nextgroupplot[width=4.6cm, xbar, bar width=6pt, xmin=0, xmax=36, ymin=0.4, ymax=6.6,
  ytick={1,...,6},
  yticklabels={\heb{שיעור המנוי},\heb{תמהיל השער},\heb{עלות \en{GPU}},\heb{נטישה חודשית},\heb{עלות הגיוס (\en{CAC})},\heb{המרה אצל השותף}},
  xlabel={\heb{אחוז מערך המידע הכולל}}, xmajorgrids, grid style={black!10},
  nodes near coords, nodes near coords style={font=\scriptsize,/pgf/number format/fixed,/pgf/number format/precision=0}]
  \addplot[fill=brandhighlight,draw=none] table[x=evppi_share_pct,y=pos]{figures/data/voi_ranking.csv};
\end{groupplot}
\end{tikzpicture}
\end{reportfigure}

\textbf{ערך המידע: מה כדאי למדוד קודם.} ההחלטה שעל הפרק היא להמשיך לפיילוט או לעצור, מול הסף \num{3.8}. ערך המידע המושלם על כל ההנחות יחד (\en{EVPI}) הוא \num{\$0.68} למשתמשת מגויסת, כשליש מעלות הגיוס שלה. הפירוק לפי מדידה (איור~\ref{fig:voi}, מימין) הופך את האינטואיציה. ניסוי ה-\en{A/B} אצל השותף, שמודד המרה ותוספת המרה, שווה \num{26\%} מערך המידע; פיילוט שמודד את עלות הגיוס שווה \num{12\%}; מעקב נטישה שווה \num{3\%}. מדידה מדויקת יותר של עלות ה-\en{GPU}, הדאגה הראשונה של כל פרויקט גנרטיבי, אינה שווה דבר: בשום נקודה בטווח שלה ההחלטה אינה מתהפכת. לכן הפיילוט מודד קודם את ההמרה אצל השותף (\en{RSK-BIZ-04}), ורק בסוף את עלויות החישוב.

\section{מחסומי כניסה, בידול ופלחי לקוחות}
\label{sec:moat}

שוק שבו פועלות \en{L'Oréal} ו-\en{Walmart} מעלה שאלה מתבקשת: מה ימנע מהן להוסיף את היכולות האלה? אף מחסום אינו מוחלט. אמידת תאורה מהעין מוכרת בספרות \cite{nishino2004eyes}; החידוש הוא בשילוב עם ההבזק וההתפלגות הקודמת, ולכן זהו יתרון זמן ולא פטנט. \textbf{נתונים:} אין מאגר ציבורי מזווג של סלפי וספקטרופוטומטר, והוא נבנה בניסוי האימות. \textbf{אינטגרציה:} המתחרים בנויים סביב קטגוריה אחת, ואילו הצבע המכויל ווקטור המידות משרתים את כולן. \textbf{כלכלה:} הדמיה חינמית בלי מימון נושאת עלות שגדלה עם השימוש, ומנוע השכבות הופך אותה למרכז רווח.

\textbf{אסטרטגיית הכניסה לשוק.} הכניסה מתוכננת בשלושה צעדים, שכל אחד מותנה בהוכחה של קודמו: (\num{1}) קמעונאי שותף אחד ותוסף בדף המוצר, עם ניסוי \en{A/B} שמודד החזרות והוספה לסל מול קבוצת ביקורת; (\num{2}) האפליקציה הנלווית, בהשקה חברתית בשני גלים; (\num{3}) רישוי לקמעונאיות נוספות רק אחרי ירידה מוכחת בהחזרות.

\textbf{אות ביקוש אורגני, עוד לפני ההשקה.} סרטונים קצרים שפורסמו על התוסף ב-\en{TikTok} וב-\en{Instagram Reels} הגיעו, בלי קידום ממומן, למאות אלפי צפיות, לאלפי סימוני אהבתי ולמאות תגובות ששאלו מה שם האפליקציה ואיך מקבלים גישה. אלה סדרי גודל כפי שנצפו בחשבונות הצוות, ולא מדידה מבוקרת: צפייה אינה משתמשת, ותגובה היא כוונה ולא התקנה. ובכל זאת זהו הנתון הראשון בדוח שמקורו בקהל. המשפך שהוא מתווה, בערך תגובת ביקוש אחת לאלף צפיות, חסר רק את החוליה האחרונה: כמה מהשואלות מתקינות. זו עלות הגיוס, השנייה בדירוג ערך המידע (איור~\ref{fig:voi}), ולכן לפני ההשקה ייפתח קישור הרשמה מוקדמת עם מעקב, שיהפוך תגובות לספירה.

\textbf{השקה חברתית בשני גלים.} \textbf{גל ראשון, מנוי \en{Pro} תמורת תוכן:} יוצרות ויוצרי תוכן קטנים מקבלים מנוי \en{Pro} (\num{\$7.99} לחודש) ומפרסמים תוכן על האפליקציה, בלי הוצאת מזומן; המנוי מחזיר את עצמו כבר במשתמשת הפעילה השלישית שיוצרת מביאה. ויש יתרון סטטיסטי: עשרות יוצרות קטנות הן עשרות מדידות בלתי תלויות של שיעור ההמרה, כפי שמשפחות תאורה רבות מלמדות יותר ממשתתפות רבות (סעיף~\ref{sec:posterior}). \textbf{גל שני, שלוש יוצרות בתשלום:} חלק מתקציב הפרסום מוקצה לשלוש יוצרות בעלות חשיפה גבוהה במיוחד, כדי להגיע לקהל היעד מהר ובהיקף. הסיכון מרוכז בשלוש יחידות, ולכן הגל השני מותנה בתוצאות הראשון, וכלל ההחלטה נגזר מכלכלת היחידה: כדי ש-\en{LTV/CAC} יישאר מעל \num{3.8}, עלות הגיוס אינה יכולה לעלות על \num{\$3.17}, כלומר כל \num{\$1,000} שמשולמים ליוצרת חייבים להביא לפחות \num{316} משתמשות פעילות. אם שיעורי ההמרה שיימדדו בגל הראשון אינם תומכים ביחס הזה, התקציב נשאר בגל הראשון.

\section{קהל היעד ומסע המשתמש}
\label{sec:audience}

הקהל של \en{AURAM} רחב וחוצה גילים ומגדרים, אבל כל פלח מגיע מסיבה אחרת ומשלם בערוץ אחר (טבלה~\ref{tab:segments}). כל ערוץ כזה כבר נמצא בכלכלת היחידה, כך שהקהל אינו הנחה נפרדת.

\begin{reporttable*}[\footnotesize]{פלחי הקהל: רגע השימוש, הצעת הערך וערוץ ההכנסה (הערכת הצוות)}{tab:segments}{P{2.6cm} P{3.5cm} R P{4.1cm}}
\headrow \hc{פלח} & \hc{רגע השימוש} & \hc{הצעת הערך} & \hc{ערוץ הכנסה} \\ \midrule
נערות & גלילה ברשת, קנייה לפי טרנד & טרנד שמתאים לגוון ולמבנה הגוף שלי & חסות מותג, מיקרו-אינטראקציה \\
נשים בכל גיל & בוקר מול הארון, קנייה מקוונת & לבוש, איפור וציפורניים במכלול אחד & שותפים, מנוי פרימיום \\
גברים מודעי-טיפוח & קנייה מקוונת, אירוע מיוחד & מידה וגוון נכונים בלי מומחיות & שותפים, מנוי פרימיום \\
קמעונאיות ביוטי ואופנה & דף המוצר באתר & פחות החזרות, יותר המרות & רישוי לפי הזמנה \\
מותגים קטנים & קטלוג עונתי & כלי שהיה זמין רק לקבוצות גדולות & תשלום לפי \en{API} \\
\end{reporttable*}

אותה תשתית נמכרת בכמה ערוצים, וזו פעולת מנע לסיכון \en{RSK-BIZ-04}. מסע המשתמש זהה לכל הפלחים (איור~\ref{fig:journey}): ארבעה צעדים, כל אחד עם תקציב זמן והדרישה שמודדת אותו.

% one-off (v3.6): the cheap \needspace breaks before this figure and the box left 38% of a page blank
{\renewcommand{\needspace}[1]{}%
\begin{reportfigure*}{מסע המשתמש בארבעה צעדים, מהרישום ועד הקנייה, עם תקציב הזמן והדרישה של כל צעד}{fig:journey}
\begin{tikzpicture}[x=1cm,y=1cm,
  st/.style={draw=#1,fill=brandcream!60!white,rounded corners=5pt,line width=1pt,text width=2.62cm,
             minimum height=2.25cm,align=center,font=\scriptsize},
  num/.style={circle,fill=#1,text=brandcream,font=\small\bfseries,inner sep=1.5pt,minimum size=0.55cm}]
  \node[st=brandprimary] (s1) at (9.55,0) {\textbf{\heb{רישום ב-\num{30} שניות}}\\[2pt]\heb{צבע עיניים, גובה, משקל,}\\\heb{היקפים; \num{3} צילומים:}\\\heb{פנים, גוף, ידיים}\\[2pt]\mutedtext{\en{REQ-VIS-02}}};
  \node[st=brandsecondary] (s2) at (6.37,0) {\textbf{\heb{עוזרת יומית וקנייה}}\\[2pt]\heb{לבוש, גוון איפור ולק}\\\heb{לפי טרנדים ותורת}\\\heb{הצבע העונתית}\\[2pt]{\tiny\mutedtext{\en{REQ-LAT-05, REC-10}}}};
  \node[st=brandaccent!85!black] (s3) at (3.18,0) {\textbf{\heb{הדמיה בשלוש שכבות}}\\[2pt]\heb{צללית מיידית; הדמיה}\\\heb{ממומנת, אינטראקציה של}\\\heb{\num{3} ש׳ או מנוי \en{VIP}}\\[2pt]{\tiny\mutedtext{\en{REQ-GEN-06, BIZ-12}}}};
  \node[st=brandhighlight] (s4) at (0,0) {\textbf{\heb{קנייה בלחיצה}}\\[2pt]\heb{כרטיסי פריטים}\\\heb{ממומנים אצל}\\\heb{קמעונאי שותף}\\[2pt]\mutedtext{\en{REQ-REC-10}}};
  \foreach \n/\c/\t in {s1/brandprimary/1,s2/brandsecondary/2,s3/brandaccent!85!black/3,s4/brandhighlight/4}
    \node[num=\c] at (\n.north) {\t};
  \draw[-Stealth,very thick,brandaccent] (s1) -- (s2); \draw[-Stealth,very thick,brandaccent] (s2) -- (s3);
  \draw[-Stealth,very thick,brandaccent] (s3) -- (s4);
  \node[font=\scriptsize,text=brandsecondary] at (9.55,-1.68) {$\le$ 30 s};
  \node[font=\scriptsize,text=brandsecondary] at (6.37,-1.68) {$\le$ 1.8 s};
  \node[font=\scriptsize,text=brandsecondary] at (3.18,-1.68) {$\le$ 3 s + $\le$ 4.5 s};
  \node[font=\scriptsize,text=brandsecondary] at (0,-1.68) {\heb{לחיצה אחת}};
\end{tikzpicture}
\end{reportfigure*}

החיכוך מרוכז בצעד הראשון, ולכן הוא מוגבל לחצי דקה ומתבצע פעם אחת; ההמלצה והצללית המיידית לעולם אינן נעולות מאחורי שער.

\begin{examplebox}[מקור הרעיון]
הרעיון נולד מול מראה בחדר אמבטיה עם תאורה גרועה, ומחבילה שנייה שחזרה לשולח בגלל מידה. \textbf{העין כסרגל צבע:} מנקודת מבט סטטיסטית, צבע העיניים המדווח אינו שדה בטופס אלא התפלגות קודמת חינמית, שהספרות בראייה ממוחשבת לא ניצלה. \textbf{כלכלה שמממנת את עצמה:} הבעיה של הדמיה גנרטיבית אינה האיכות אלא מי משלם על ה-\en{GPU}, ולכן המימון הפך לרכיב בארכיטקטורה. \textbf{צורך יומיומי:} נערה לפני מסיבה, אישה לפני פגישה וגבר שמקפיד על טיפוח שואלים אותה שאלה: "זה באמת יתאים לי?".
\end{examplebox}}
